%% --------------------------------------------------------------------------
\section{Structural and Property Microbenchmarks}
\label{sec:eval:microbenchmarks}

Composite graph queries interleave structural traversal with property access, so no single layer accounts for the measured latency.
We first isolate per-operation costs using microbenchmarks from the AsterDB artifact (structural mutations, neighbor lookups, and property CRUD) before evaluating composite queries in~\autoref{sec:eval:property-ldbc}.
Structural operations are measured on two undirected graphs: DBLP co-authorship (${\sim}$317\,K vertices, ${\sim}$1\,M edges) and WikiTalk (${\sim}$2.4\,M vertices, ${\sim}$5\,M edges).
Property operations are measured on two property graphs from the AsterDB artifact, LDBC and Freebase, with Freebase substantially larger.
The baseline systems execute each operation through an embedded TinkerPop provider loaded in the Gremlin Console;\footnote{JanusGraph is the exception: it runs a Gremlin Server in a separate process on the same host, so the console reaches it over a local socket and its measurements carry request serialization on top of the provider cost.}
\project{} bypasses this layer, accessing its storage engine through a native C++ API; AsterDB is measured through its own TinkerPop provider, which is backed by RocksDB for structural operations only.
\project{} property operations use the embedded property storage mode.
The two operation classes run under different harnesses.
Structural operations run 100{,}000 iterations inside a single container, so the reported cost reflects a warmed JVM.
Property operations run once per container for the baseline systems, and 10 times in a single process for \project{} and AsterDB.
We report the mean latency.

%%%%%%%%%%%%%%%%%%%%%%%%%%%%%%%%%%%%%%%%%%%%%%%%%%%%%%%%%%%%%%%%%%%%%%%%%%%%%%%%
\begin{table}[t]
  \centering
  % \footnotesize
  \setlength{\tabcolsep}{2pt}
  \begin{tabular*}{\columnwidth}{@{\extracolsep{\fill}}ll rrrrrr@{}}
    \toprule
    Operation & Graph & Neo4j & ArangoDB & PostgreSQL & AsterDB & FG:Adj & FG:Ekey \\
    \midrule
    \multirow{2}{*}{Add Vertex}
    & DBLP & 8.75 & 14{,}065 & 87 & 11.89 & \textbf{1.96} & 1.97 \\
    & Wiki & 7.75 & 70{,}248 & 84 & 11.77 & 2.46 & \textbf{2.38} \\
    \midrule
    \multirow{2}{*}{Add Edge}
    & DBLP & 16.55 & 9{,}262 & 160 & 16.43 & 11.35 & \textbf{7.01} \\
    & Wiki & \textbf{20.35} & 10{,}645 & 186 & 38.54 & 49.42 & 78.23 \\
    \midrule
    \multirow{2}{*}{Del Edge}
    & DBLP & 22.09 & 17{,}864 & 60 & 13.32 & 10.91 & \textbf{6.66} \\
    & Wiki & \textbf{24.34} & 21{,}506 & 66 & 26.84 & 24.86 & 44.46 \\
    \midrule
    \multirow{2}{*}{Get Nbrs}
    & DBLP & 11.67 & 43{,}104 & 32{,}205 & 31.42 & \textbf{1.13} & 1.85 \\
    & Wiki & 50.61 & 37{,}146 & 2{,}581 & 49.60 & \textbf{2.10} & 5.17 \\
    \bottomrule
  \end{tabular*}
  \caption{Structural operation latency (\textmu{}s) on DBLP and WikiTalk.
    OrientDB, JanusGraph, and NebulaGraph excluded (no WikiTalk results within the 2-hour timeout).
    The best (lowest) latency per row appears in \textbf{bold}.}
  \label{tab:structural-ops}
\end{table}
%%%%%%%%%%%%%%%%%%%%%%%%%%%%%%%%%%%%%%%%%%%%%%%%%%%%%%%%%%%%%%%%%%%%%%%%%%%%%%%%

\textbf{Structural operations.}
\autoref{tab:structural-ops} and \autoref{fig:structural-ops} report the latency of four structural primitives.
\project{} Adjacency~List is fastest for neighbor lookups: 1.13\,$\mu$s on DBLP and 2.10\,$\mu$s on WikiTalk.
It is $10\times$ faster than Neo4j on DBLP (11.67\,$\mu$s) and $24\times$ faster on WikiTalk (50.61\,$\mu$s), and four orders of magnitude faster than ArangoDB.
A neighbor lookup in \project{} translates to a single B-tree cursor seek to the source vertex's adjacency list, followed by a sequential scan of packed neighbor entries within the same page.
Neo4j follows pointer chains in its native store files, incurring a seek per hop; ArangoDB resolves edges through document-level index lookups, paying per-edge overhead.

Vertex insertion completes in under 2.5\,$\mu$s for \project{} regardless of graph size, $4\times$ faster than Neo4j (8--9\,$\mu$s).\footnote{This comparison is conservative for \project{}: Neo4j's structural harness commits every 10{,}000 operations, amortizing per-transaction overhead across the batch, while each \project{} mutation wraps its own \wt{} transaction.}
Edge insertion reveals a degree-dependent tradeoff.
On DBLP, EdgeKey inserts an edge in 7.01\,$\mu$s, faster than Neo4j (16.55\,$\mu$s), while Adjacency~List takes 11.35\,$\mu$s.
On the denser WikiTalk graph, Neo4j (20.35\,$\mu$s) outperforms both \project{} layouts (Adjacency~List: 49.42\,$\mu$s, EdgeKey: 78.23\,$\mu$s).
Under Adjacency~List, each edge insert appends to the source vertex's packed neighbor list, a read-modify-write on the B-tree row whose cost grows with vertex degree.
EdgeKey stores each edge as an independent row, avoiding this cost, but \wt{}'s per-transaction overhead (begin, cursor insert, commit) is the bottleneck at this granularity.
Edge deletion follows a similar pattern: EdgeKey is fastest on DBLP (6.66\,$\mu$s vs.\ Neo4j's 22.09\,$\mu$s), while on WikiTalk, Neo4j again outperforms both layouts (24.34\,$\mu$s vs.\ Adjacency~List's 24.86\,$\mu$s and EdgeKey's 44.46\,$\mu$s).
The degree-dependent cost of Adjacency~List mutations and the per-transaction overhead of EdgeKey both grow relative to Neo4j's pointer-based store as graphs become denser.

ArangoDB is three to four orders of magnitude slower than \project{} on every structural operation.
PostgreSQL is $3$--$44\times$ slower on mutations and up to $28{,}000\times$ slower on neighbor lookups.
ArangoDB stores edges as JSON documents in a separate collection, requiring document-level index lookups and writes for each mutation.
PostgreSQL, accessed through the SQLG TinkerPop provider, translates each graph operation into SQL statements on relational tables, paying relational join and transaction overhead.
A direct-SQL decomposition analogous to \autoref{tab:crud-overhead} would isolate the SQLG translation cost, but the benchmark artifact does not expose a native SQL path for PostgreSQL.

%%%%%%%%%%%%%%%%%%%%%%%%%%%%%%%%%%%%%%%%%%%%%%%%%%%%%%%%%%%%%%%%%%%%%%%%%%%%%%%%
\begin{figure}[h]
  \centering
  \includegraphics[width=0.75\columnwidth]{property/structural_ops.pdf}
  \caption{Structural operation latency (\textmu{}s, log scale) on DBLP and WikiTalk.
    ArangoDB bars extend to $10^4$--$10^5$\,\textmu{}s.}
  \label{fig:structural-ops}
\end{figure}
%%%%%%%%%%%%%%%%%%%%%%%%%%%%%%%%%%%%%%%%%%%%%%%%%%%%%%%%%%%%%%%%%%%%%%%%%%%%%%%%

\textbf{Property mutations.}
\autoref{fig:property-crud} reports eight property CRUD operations across two datasets.
On the LDBC dataset, \project{} completes a vertex property update in ${\sim}$10\,$\mu$s, against 47\,ms for Neo4j, 24\,ms for ArangoDB, and 106\,ms for JanusGraph.
These gaps of three to four orders of magnitude measure the artifact's TinkerPop path rather than the baselines' storage engines.
The decomposition that follows separates the two.

The artifact runs each baseline property operation in a fresh container: a new JVM opens the database from disk and resolves the target element, then times a single property mutation and transaction commit before exiting.
The timed path therefore executes without JIT-compiled code and with a cold storage-engine cache, both consequences of the single-operation-per-container harness rather than of TinkerPop itself.
The structural harness amortizes these costs across a single container of 100{,}000 operations, which is why Neo4j appears in the microsecond range in \autoref{tab:structural-ops} but in the tens of milliseconds here.

AsterDB completes property mutations in 40--110\,$\mu$s, but its Gremlin bridge writes properties to TinkerGraph's in-memory map, so they never reach its RocksDB backend.
\project{} issues a single \wt{} cursor write with an implicit commit, positioning by key and writing the property value into the B-tree; it persists at checkpoint, so neither system pays a per-operation fsync here.
AsterDB's figures therefore isolate the TinkerPop provider path in a warm JVM with the storage term removed, and bound provider overhead at $4$--$10\times$ the cost of a \project{} mutation that carries full B-tree and transaction work.

Among the baseline systems, OrientDB is consistently fastest on mutations (8--18\,ms), suggesting a lightweight provider implementation.
JanusGraph is slowest (21--207\,ms): it carries request serialization to its Gremlin Server process on top of a layered architecture (TinkerPop provider, JanusGraph core, pluggable storage backend) that compounds on fine-grained mutations.

\newflag[new: CRUD overhead decomposition]
\textbf{Overhead decomposition.}
We measure the same six CRUD operations on Neo4j and ArangoDB through their native query languages (Cypher and AQL, respectively), bypassing the TinkerPop layer entirely.
Both systems run as warm servers with many operations per session, so this measurement also removes the cold-start costs described above.
Each operation runs in two modes: single-shot (one operation per client request) and batched (1{,}000 operations per request, amortizing query parsing, planning, and transaction commit).
\autoref{tab:crud-overhead} reports all six mutations for both systems.
On Neo4j, measuring the same mutation on the warm native path drops vertex property update latency from 46.68\,ms to 3.27\,ms, a $14\times$ reduction that accounts for 93\% of the total gap.
Batching further amortizes per-request overhead (query parse and plan, transaction begin and commit, client round-trip), reducing per-operation cost to 0.027\,ms; \project{} reaches 0.011\,ms, $2.5\times$ faster.
Across both systems, the TinkerPop path and its cold-start costs together account for 92--96\% of the gap on insert and update operations.
Neo4j's deletes show a smaller share (20--68\%), because its TinkerPop delete path already runs faster (3.7--7.2\,ms), leaving per-request overhead to dominate the remainder; ArangoDB's deletes retain a high share (77--90\%).
Through batched Cypher, Neo4j's storage engine is $2.5$--$5\times$ slower than \project{}, not the $4{,}000\times$ that the artifact's TinkerPop path reports.
ArangoDB shows the same pattern: native AQL reduces vertex property update latency from 23.86\,ms to 1.64\,ms ($15\times$), with the TinkerPop path accounting for 93\% of the gap.
ArangoDB's single-shot native latency (1.5--1.8\,ms) falls below Neo4j's (2.3--3.3\,ms), indicating a more efficient per-request path.
Batched vertex operations are comparable across the two systems (${\sim}$0.02\,ms/op), while ArangoDB's batched edge operations cost $3$--$4\times$ more (0.14 vs.\ 0.04\,ms/op), reflecting its document-oriented edge storage.

For OrientDB and PostgreSQL, the benchmark artifact does not expose a native-API path, so we cannot perform the same decomposition; JanusGraph's native API is TinkerPop itself, leaving no layer to remove.
The decomposition establishes the engine-level comparison for both systems: under batched native queries, \project{} is $2.5$--$5\times$ faster than Neo4j, and $2$--$3\times$ faster than ArangoDB on vertex properties and $11$--$16\times$ on edge properties.
Subsequent claims about engine-level performance rest on this comparison.
\autoref{fig:property-crud} reports the artifact's single-operation-per-container path, so its baseline figures represent an upper bound on each system's engine-level gap rather than a direct measurement of it.
\autoref{tab:structural-ops} is measured on the warm path and needs no such correction, though its baseline measurements commit every 10{,}000 operations against a per-operation transaction for \project{}, understating \project{}'s advantage.
Both subsequent experiments measure the baselines on a native path.
The LDBC query workload (\autoref{sec:eval:property-ldbc}) executes Neo4j and ArangoDB through their native query languages with per-request overhead amortized by batching, and the graph algorithm comparison (\autoref{sec:eval:property-graphalytics}) quantifies the TinkerPop overhead and reports the residual gap.
%%%%%%%%%%%%%%%%%%%%%%%%%%%%%%%%%%%%%%%%%%%%%%%%%%%%%%%%%%%%%%%%%%%%%%%%%%%%%%%%
\begin{table}[t]
  \centering
  \setlength{\tabcolsep}{2pt}
  \begin{tabular*}{\columnwidth}{@{\extracolsep{\fill}}l rrrr rrr@{}}
    \toprule
    & \multicolumn{4}{c}{Latency (ms)} & \multicolumn{3}{c}{Overhead (\%)} \\
    \cmidrule(lr){2-5} \cmidrule(lr){6-8}
      Operation & TP & Native & Batch/op & FG & TP & Native & Engine \\
    \midrule
    \multicolumn{8}{l}{\textit{Neo4j (Cypher)}} \\
    Update vprop & 46.68 & 3.27 & 0.027 & 0.011 & 93.0 & 6.9 & 0.0 \\
    Insert vprop & 68.44 & 2.94 & 0.026 & 0.008 & 95.7 & 4.3 & 0.0 \\
    Delete vprop &  3.71 & 2.96 & 0.028 & 0.007 & 20.3 & 79.2 & 0.6 \\
    Update eprop & 41.96 & 2.33 & 0.043 & 0.013 & 94.5 & 5.5 & 0.1 \\
    Insert eprop & 42.66 & 2.65 & 0.041 & 0.010 & 93.8 & 6.1 & 0.1 \\
    Delete eprop &  7.24 & 2.30 & 0.044 & 0.009 & 68.3 & 31.2 & 0.5 \\
    \midrule
    \multicolumn{8}{l}{\textit{ArangoDB (AQL)}} \\
    Update vprop & 23.86 & 1.64 & 0.023 & 0.011 & 93.2 & 6.8 & 0.0 \\
    Insert vprop & 22.48 & 1.76 & 0.020 & 0.008 & 92.2 & 7.7 & 0.1 \\
    Delete vprop &  7.13 & 1.65 & 0.019 & 0.007 & 76.9 & 22.9 & 0.2 \\
    Update eprop & 20.97 & 1.72 & 0.146 & 0.013 & 91.9 & 7.5 & 0.6 \\
    Insert eprop & 23.42 & 1.64 & 0.143 & 0.010 & 93.1 & 6.4 & 0.6 \\
    Delete eprop & 14.34 & 1.51 & 0.142 & 0.009 & 89.5 & 9.6 & 0.9 \\
    \bottomrule
  \end{tabular*}
  \caption{CRUD overhead decomposition on LDBC for Neo4j and ArangoDB.
    TP = the artifact's TinkerPop measurement, one operation per fresh container, and so inclusive of cold-start costs; Native = single-shot native query (Cypher, AQL) against a warm server; Batch/op = 1{,}000-op batch, per-operation; FG = \project{} Adjacency~List.
    The TP, Native, and Engine columns show each layer's share of the total gap (TP $-$ FG) as a percentage.}
  \label{tab:crud-overhead}
\end{table}
%%%%%%%%%%%%%%%%%%%%%%%%%%%%%%%%%%%%%%%%%%%%%%%%%%%%%%%%%%%%%%%%%%%%%%%%%%%%%%%%
%%%%%%%%%%%%%%%%%%%%%%%%%%%%%%%%%%%%%%%%%%%%%%%%%%%%%%%%%%%%%%%%%%%%%%%%%%%%%%%%
\begin{figure}[h!]
  \centering
  \includegraphics[width=0.85\columnwidth]{property/property_crud.pdf}
  \caption{Property CRUD latency (ms, log scale) on LDBC and Freebase.
    Baseline measurements run one operation per fresh container and so include cold-start costs that the harness never amortizes; AsterDB's property writes go to an in-memory map and never reach its storage engine.
    \autoref{tab:crud-overhead} separates these costs from engine performance.
    Property search narrows the gap as scan work outweighs overhead.}
  \label{fig:property-crud}
\end{figure}
%%%%%%%%%%%%%%%%%%%%%%%%%%%%%%%%%%%%%%%%%%%%%%%%%%%%%%%%%%%%%%%%%%%%%%%%%%%%%%%%

\textbf{Property search.}
Property search scans all vertices or edges to find entries matching a specific property value, analogous to an unindexed \code{WHERE} clause.
The performance gap narrows here, because the scan work outweighs per-query overhead.
On LDBC, \project{} Adjacency~List completes a vertex property search in 14.4\,ms, $29\times$ faster than Neo4j (425\,ms) and $188\times$ faster than ArangoDB (2{,}712\,ms).
On the larger Freebase graph, the same operation takes 2.1\,s versus Neo4j's 27.3\,s ($13\times$).

EdgeKey is slower than Adjacency~List on property search: 153\,ms versus 14.4\,ms on LDBC vertex search ($10.6\times$).
Under EdgeKey, each edge occupies its own B-tree row, so a full scan touches many more B-tree pages than does Adjacency~List, where edges are packed into per-vertex records.
The denser packing gives better cache and I/O efficiency for sequential scans.

ArangoDB exhibits an extreme outlier on LDBC edge property search: ${\sim}$1{,}904\,s (${\sim}$32\,minutes), likely caused by a missing secondary index forcing a cross-collection join over the entire edge set.